\documentclass[10pt]{article}

% ---------- Document Identity (update these only) ----------
\newcommand{\DocStage}{}
\newcommand{\DocTitle}{Your Road to Tier One \textbf{SOF}}
\newcommand{\ProgramName}{182-Week Civilian-to-Selection Readiness Pipeline}
\newcommand{\ProgramSubtitle}{}

% ---------- Page ----------
\usepackage[legalpaper,left=0.5in,right=0.5in,top=0.6in,bottom=0.45in]{geometry}

% ---------- Typography ----------
\usepackage[T1]{fontenc}
\usepackage[utf8]{inputenc}
\usepackage{libertinus}
\usepackage{microtype}
\usepackage{parskip}
\usepackage{ragged2e}
\usepackage{textcomp}
\AtBeginDocument{\color{black}}
\AtBeginDocument{\pagecolor{white}}
\AtBeginDocument{\justifying}

% ---------- Landscape pages ----------
\usepackage{pdflscape}

% ---------- Color ----------
\usepackage[table]{xcolor}
\definecolor{StageBlue}{HTML}{1F4E79}
\definecolor{StageBlueDark}{HTML}{163A5A}
\definecolor{LightGray}{HTML}{F2F4F7}
\definecolor{MidGray}{HTML}{D9DEE5}
\definecolor{RuleGray}{HTML}{C7CED8}
\definecolor{HeaderInk}{HTML}{000000}
\definecolor{Ink}{HTML}{000000}

% ---------- Utility ----------
\usepackage{enumitem}
\sloppy
\setlength{\emergencystretch}{2em}

% ---------- Boxes / UI ----------
\usepackage[most]{tcolorbox}

\tcbset{
  charterTitle/.style={
    enhanced,
    colback=StageBlue!4,
    colframe=StageBlue,
    boxrule=1.25pt,
    arc=3pt,
    left=16pt,right=16pt,top=14pt,bottom=14pt,
    borderline north={3.2pt}{0pt}{StageBlue},
    borderline west={4pt}{0pt}{StageBlue},
    borderline south={1.25pt}{0pt}{StageBlue},
    drop shadow={black!25!white},
  },
  charterRule/.style={
    enhanced,
    colback=LightGray,
    colframe=RuleGray,
    boxrule=0.85pt,
    arc=3pt,
    left=12pt,right=12pt,top=10pt,bottom=10pt,
    borderline west={3pt}{0pt}{StageBlue},
  },
  charterBadge/.style={
    enhanced,
    colback=StageBlue,
    coltext=white,
    colframe=StageBlueDark,
    boxrule=0.6pt,
    arc=2pt,
    left=6pt,right=6pt,top=3pt,bottom=3pt,
  }
}
\newtcbox{\Badge}{nobeforeafter,charterBadge}

% ---------- Lists ----------
\setlist[itemize]{leftmargin=1.5em, itemsep=0.25em, topsep=0.25em}
\setlist[enumerate]{leftmargin=1.8em, itemsep=0.25em, topsep=0.25em}

% ---------- TOC ----------
\usepackage{tocloft}
\setcounter{tocdepth}{3}
\setlength{\cftbeforesecskip}{3pt}
\setlength{\cftbeforesubsecskip}{2pt}
\setlength{\cftbeforesubsubsecskip}{0pt}
\renewcommand{\cftsecfont}{\bfseries}
\renewcommand{\cftsecpagefont}{\bfseries}
\renewcommand{\cftsubsecfont}{\normalfont}
\renewcommand{\cftsubsecpagefont}{\normalfont}
\renewcommand{\cftsubsubsecfont}{\normalfont}
\renewcommand{\cftsubsubsecpagefont}{\normalfont}
\setlength{\cftsecindent}{0pt}
\setlength{\cftsubsecindent}{1.25em}
\setlength{\cftsubsubsecindent}{2.8em}
\setlength{\cftsecnumwidth}{2.1em}
\setlength{\cftsubsecnumwidth}{2.8em}
\setlength{\cftsubsubsecnumwidth}{3.6em}

% Remove the built-in TOC heading so only the custom heading is shown
\makeatletter
\renewcommand{\tableofcontents}{\@starttoc{toc}}
\makeatother

% ---------- Header/Footer: page number only ----------
\usepackage{fancyhdr}
\pagestyle{fancy}
\fancyhf{}
\renewcommand{\headrulewidth}{0pt}
\renewcommand{\footrulewidth}{0pt}
\fancyfoot[C]{\thepage}

% ---------- Section formatting ----------
\usepackage{titlesec}

\titleformat{\section}
  {\Large\bfseries\color{Ink}}
  {\thesection}{0.6em}{}
  [\vspace{0.25em}\textcolor{StageBlue}{\rule{\linewidth}{0.8pt}}]

\titleformat{\subsection}
  {\large\bfseries\color{Ink}}
  {\thesubsection}{0.6em}{}

\titleformat{\subsubsection}
  {\normalsize\bfseries\color{Ink}}
  {\thesubsubsection}{0.6em}{}

\titlespacing*{\section}{0pt}{0.95em}{0.35em}
\titlespacing*{\subsection}{0pt}{0.65em}{0.25em}
\titlespacing*{\subsubsection}{0pt}{0.55em}{0.2em}

% ---------- Table engine ----------
\usepackage{tabularray}
\UseTblrLibrary{booktabs}

% ---------- tabularray: GLOBAL table treatment ----------
\SetTblrInner{
  cells = {fg=black, bg=white, valign=t},
}

% ---------- Header helpers ----------
\newcommand{\Hdr}[1]{\shortstack[c]{\strut #1\strut}}

% ---------- Lines ----------
\newcommand{\TitleLine}{\textcolor{StageBlue}{\rule{0.98\linewidth}{0.95pt}}}
\newcommand{\SubTitleLine}{\textcolor{RuleGray}{\rule{0.98\linewidth}{0.65pt}}}

% ---------- Continued on next page ----------
\newcommand{\ContinuedOnNextPageInline}{%
  \par
  \vfill
  \noindent\textcolor{RuleGray}{\rule{\linewidth}{1pt}}\vspace{-2pt}
  \noindent\hfill\textit{Continued on next page}\par\vspace{-10pt}
  \noindent\textcolor{RuleGray}{\rule{\linewidth}{1pt}}
  \clearpage
}

% ---------- PDF hyperlinks ----------
\usepackage[hidelinks]{hyperref}
\hypersetup{
  colorlinks=false,
  pdfborder={0 0 0},
  linktoc=all,
  pdftitle={\DocTitle},
  pdfauthor={\ProgramName}
}

% =========================================================
\begin{document}

\begin{tcolorbox}
\begin{center}
Stage 11 — Pre-Selection Conditioning and Recovery\\
\textbf{SOF} Physical Development Transformation Program\\
182-Week Civilian-to-Selection Readiness Pipeline
\TitleLine
\end{center}
\end{tcolorbox}

\vspace{0.35em}

\section*{Stage 11 — Pre-Selection Conditioning and Recovery}
\phantomsection
\addcontentsline{toc}{section}{Stage 11 — Pre-Selection Conditioning and Recovery}

\subsection*{Introduction}
\phantomsection
\addcontentsline{toc}{subsection}{Introduction}

\subsubsection*{1) Purpose}

Stage 11 exists to finalize execution doctrine for the selection approach period across the final phases, with explicit focus on maintaining admissible evidence, controlling risk, and preventing last-minute drift. Stage 11 provides operational governance and implementation posture, not programming, and it does not redesign phase missions, tests, thresholds, calendars, or equipment timing.

Stage 11 leads into the locked deliverable set:

\begin{itemize}
\item 11.A — Pre-Selection Conditioning Overview
\item 11.B — Final Preparation and Test Optimization
\item 11.C — Recovery and Injury Prevention in Final Phases
\item 11.D — Nutritional Adjustments for Final Selection
\end{itemize}

\subsubsection*{2) Scope}

\paragraph*{Inclusions}

\begin{itemize}
\item 11.A: pre-selection conditioning posture and constraints at governance level (no workouts; no programming detail).
\item 11.B: final preparation and test optimization posture that protects admissibility, preserves comparability, and avoids novelty in protected windows.
\item 11.C: recovery and injury-prevention overlays for final phases, framed non-medically and aligned to stop-rule authority and \textbf{MEDICAL} \textbf{HOLD} posture where indicated.
\item 11.D: nutrition adjustments for final selection that implement Stage 10 (10.A–10.D) without redefining it, including hydration/electrolyte stability consistent with Stage 10.C.
\item Validity protection:
\begin{itemize}
\item adherence to Stage 1.F validity tags and admissibility posture (\textbf{VALID} / \textbf{MODIFIED} / \textbf{INVALID} for sessions; \textbf{VALID} / \textbf{INVALID} for tests),
\item adherence to Stage 2 protocol card execution posture and Stage 2.E required fields (Environment Unit IDs, Tool Standard, Evidence Ref where applicable, missing-data doctrine).
\end{itemize}
\item Scheduling protection:
\begin{itemize}
\item adherence to Stage 3 protected Deload/Test windows,
\item reslot posture for invalid/unsafe tests,
\item no novelty near gates/test weeks.
\end{itemize}
\end{itemize}

\paragraph*{Exclusions}

\begin{itemize}
\item No new training content or programming prescriptions (no workouts, sets/reps, intervals, weekly splits, or day-by-day schedules).
\item No new tests, thresholds, domain tags, phase purposes, or equipment categories.
\item No clinical diagnosis or treatment planning.
\item No prescription guidance and no prohibited substance guidance.
\item No changes to Stage 6–7 capability timing; Stage 11 cannot require capabilities earlier than Stage 6 indicates.
\end{itemize}

\subsubsection*{3) How to Use}

\begin{itemize}
\item Consult Stage 11 during Phases 11–13 and any time the program is inside a protected Deload/Test window where evidence is collected for gates.
\item Apply Stage 1.F stop rules and \textbf{MEDICAL} \textbf{HOLD} posture without drift:
\begin{itemize}
\item safety overrides schedule and performance,
\item red flags terminate testing and trigger reslot discipline rather than forced attempts.
\end{itemize}
\item Preserve admissible evidence:
\begin{itemize}
\item execute tests only under Stage 2 protocol posture,
\item enforce Stage 2.E required fields (environment IDs, tool standard, evidence refs),
\item treat missing required fields as inadmissible and reslot.
\end{itemize}
\item Prevent last-minute novelty:
\begin{itemize}
\item do not introduce new shoes, routes, tools, hydration practices, or nutrition structure inside the no-novelty windows governed upstream,
\item if an unavoidable change occurs, log it explicitly and interpret outcomes conservatively.
\end{itemize}
\end{itemize}

\paragraph*{Trainee operational use}

\begin{itemize}
\item Stage 11 outputs are required reading for the selection approach period once issued:
\begin{itemize}
\item execute with honesty posture in logs,
\item report red flags immediately,
\item do not self-override stop rules or attempt to “prove it” on compromised days.
\end{itemize}
\item Maintain controlled variables near gates:
\begin{itemize}
\item keep nutrition/hydration routines stable (Stage 10),
\item keep equipment and environment identity stable (Stage 2.E).
\end{itemize}
\end{itemize}

\paragraph*{Decision posture}

\begin{itemize}
\item Gate outcomes remain: \textbf{ADVANCE} / \textbf{HOLD} / \textbf{REGRESS} / \textbf{MEDICAL} \textbf{HOLD}.
\item Decisions rely on admissible \textbf{VALID} evidence only; invalid tests provide no gate credit and trigger reslot per Stage 3.
\end{itemize}

\subsubsection*{4) Inputs and Assumptions}

\paragraph*{Inputs}

\begin{itemize}
\item Stage 1.F governance and validity doctrine: validity tags, admissibility, stop posture, gate outcomes, freeze windows, patch discipline.
\item Stage 2 protocol cards and \textbf{MET} inventory plus Stage 2.E required logging fields:
\begin{itemize}
\item Protocol Cards C# and Metric IDs \textbf{MET-2B}-###,
\item Environment Unit IDs and Tool Standard posture,
\item missing-data doctrine: \textbf{UNKNOWN} / \textbf{MISSING}; no inference.
\end{itemize}
\item Stage 3 protected Deload/Test windows, spacing intent, and reslot rules for invalid/unsafe tests.
\item Stage 5 Phase 11–Phase 13 intent and gate posture (Stage 11 cannot redefine phase missions or gate structure).
\item Stage 6–7 capability readiness timing and category manuals (no early requirements).
\item Stage 10 nutrition and hydration operating systems (Stage 11 applies them operationally; does not redefine them).
\end{itemize}

\paragraph*{Assumption Ledger}

\begin{itemize}
\item Access to testing environments may vary (route availability, pool availability, facility access), requiring substitutions or reslotting consistent with validity rules.
\item Selection timeline may be externally fixed, increasing the need for early verification of feasibility and strict no-stacking reslot posture.
\item Environment volatility (heat, ice, air quality, visibility) may require indoor substitutions or reslot decisions consistent with Stage 3 and Stage 2.E logging posture.
\end{itemize}

\subsubsection*{5) Interfaces}

\paragraph*{Upstream dependencies}

\begin{itemize}
\item Stage 1.F: validity, admissibility, stop posture, gate outcomes, freeze windows, patch discipline.
\item Stage 2 and Stage 2.E: protocol execution and required logging fields that make evidence admissible.
\item Stage 3: protected windows, reslot doctrine, and no novelty near gates/test weeks posture.
\item Stage 5: Phase 11–13 mission intent and gate structure.
\item Stage 6–7: capability timing and category-level operating manuals.
\item Stage 10: nutrition operating system and hydration/electrolyte rules.
\end{itemize}

\paragraph*{Downstream consumers}

\begin{itemize}
\item Stage 11 deliverables 11.A–11.D.
\item Stage 9 packaging and change control when Stage 11 is bundled into a release package; Known Issues and \textbf{ISS-08}-### references must be preserved.
\end{itemize}

\paragraph*{Crosswalk hooks}

\begin{itemize}
\item Phase IDs and gate windows (Phase 11–Phase 13 and their protected windows).
\item Stage 2 identifiers: C# and \textbf{MET-2B}-### referenced without inventing IDs.
\item Environment Unit IDs and Tool Standard posture (Route \textbf{ID} / Pool \textbf{ID} / Machine \textbf{ID} where applicable).
\item Issue IDs \textbf{ISS-08}-### where applicable for known constraints that affect feasibility or admissibility.
\item Change-control references:
\begin{itemize}
\item \textbf{PATCH} and \textbf{CHANGE-LOG} entries for any controlled updates,
\item freeze-window discipline preserved.
\end{itemize}
\end{itemize}

\subsubsection*{6) Gate / Acceptance Criteria}

Stage 11 Intro is complete when it:

\begin{itemize}
\item States Stage 11 purpose and the non-programming boundary clearly.
\item Names 11.A–11.D and describes their scope accurately without outputting their content.
\item Explicitly references:
\begin{itemize}
\item Stage 1.F validity/admissibility and gate outcomes,
\item Stage 2 protocol and Stage 2.E required fields posture,
\item Stage 3 protected windows and reslot/no-novelty posture,
\item Stage 10 nutrition/hydration application posture.
\end{itemize}
\item Preserves safety boundaries and stop-rule authority as controlling.
\item Declares crosswalk hooks sufficient for packaging, audit traceability, and issue linkage (\textbf{ISS-08}-### where relevant).
\item Flags any detected mismatch as "Requires Stage 8 review" rather than silently correcting upstream doctrine.
\end{itemize}

\subsubsection*{7) Common Failure Modes + Controls}

\begin{itemize}
\item Last-minute novelty near tests → prohibited; enforce no novelty posture and reslot if needed to preserve admissibility.
\item Chasing single-day peak performance → deny; enforce validity posture and repeatable evidence preference; reslot rather than stack.
\item Ignoring fatigue and red flags → \textbf{STOP} \textbf{NOW} posture; apply \textbf{MEDICAL} \textbf{HOLD} when indicated; do not attempt compromised tests.
\item Drifting into unlogged changes → enforce Stage 2.E required fields and Stage 1.F change-control discipline; missing data remains \textbf{UNKNOWN}/\textbf{MISSING}.
\item Equipment or environment assumptions late → verify Stage 6 timing; substitute or reslot; log explicitly; do not improvise unsafe substitutes.
\item Redefining nutrition late → implement Stage 10 only; do not invent new rules; maintain stability near protected windows.
\item Water governance drift → prohibited; no unsupervised underwater/breath-hold ever; treat any governance gap as inadmissible and reslot.
\end{itemize}

\subsubsection*{8) Versioning Notes}

\begin{itemize}
\item Stage 11 is part of Program v1.0 and is frozen unless patched under change control.
\item Any modification requires:
\begin{itemize}
\item patch record entry,
\item impact level designation,
\item revalidation of validity and timing alignment with Stage 2.E and Stage 3 protected windows.
\end{itemize}
\item Any upstream conflicts must be flagged as Requires Stage 8 review rather than silently corrected.
\end{itemize}

\end{document}